% !PRINTED 367
\weberpage{367}
If $D$ is positive, then $(E, E')$ still has to be determined.

Let
\[
\varepsilon = \frac{t + u\sqrt{D}}{2}
\]
be the fundamental unit in $O$, and $\lambda$ the smallest positive exponent for which
\[\varepsilon^{\lambda} = \frac{t + u\sqrt{D}}{2}\]
% [sic] The source leaves this second display unnumbered although the following sentence refers to (17), and reuses t,u; both features are retained.
is contained in $O'$. Then $E$ consists of the powers $\pm \varepsilon^{\nu}$, and $E'$ of the powers $\pm \varepsilon^{\lambda\nu}$, and we have
\[
(E, E') = \lambda,
\]
and formula (15) gives the class number $h'$. In general no further determination can be given for $\lambda$.

If again $\Delta \equiv 1 \, (\mathrm{mod}\ 4)$ and $Q = 2$, then $\lambda = 1$ or $= 3$; namely $= 1$ if $u$ is even in (17), and $= 3$ if $u$ is odd. The latter can occur only when $\Delta \equiv 5 \, (\mathrm{mod}\ 8)$, and consequently for $\Delta \equiv 1 \, (\mathrm{mod}\ 8)$ we always have $\lambda = 1$, while for $\Delta \equiv 1 \, (\mathrm{mod}\ 8)$ % [sic]
$\lambda$ can be either $= 1$ or $= 3$.

\begin{center}\rule{0.32\linewidth}{0.4pt}\end{center}
